% !TEX program = xelatex
% Самодостаточный конспект. Компилировать XeLaTeX/LuaLaTeX (tectonic 08-tonkie-priznaki.tex).
\documentclass[11pt,a4paper]{article}
\newcommand{\coursename}{Математический анализ}
\newcommand{\rightheader}{§8. Тонкие признаки сходимости}
\newcommand{\doctitle}{Тонкие признаки сходимости}
% ==== Общая преамбула конспектов (собирается в каждый .tex как самодостаточная) ====
% Перед этим файлом должны быть определены: \documentclass и макросы
%   \coursename  — название курса (левый колонтитул, подзаголовок),
%   \rightheader — правый колонтитул (напр. «§2. Рациональные дроби»),
%   \doctitle    — заголовок раздела на титуле.
\usepackage{fontspec}
\setmainfont{cmunrm.otf}[
  BoldFont       = cmunbx.otf,
  ItalicFont     = cmunti.otf,
  BoldItalicFont = cmunbi.otf]
\usepackage{polyglossia}
\setdefaultlanguage{russian}
\setotherlanguage{english}

\usepackage{amsmath,amssymb,amsthm,mathtools}
\usepackage[a4paper,margin=2.2cm]{geometry}
\usepackage{enumitem}
\usepackage{graphicx}
\usepackage{array}
\usepackage{xcolor}
\definecolor{accent}{HTML}{1F6FEB}
\definecolor{rulecol}{HTML}{D0D7DE}
\usepackage[colorlinks=true,linkcolor=accent,urlcolor=accent,citecolor=accent]{hyperref}
\usepackage{titlesec}
\usepackage{fancyhdr}

% ----- Колонтитулы -----
\pagestyle{fancy}
\fancyhf{}
\fancyhead[L]{\small\itshape \coursename}
\fancyhead[R]{\small\itshape \rightheader}
\fancyfoot[C]{\thepage}
\renewcommand{\headrulewidth}{0.4pt}
\renewcommand{\headrule}{\hbox to\headwidth{\color{rulecol}\leaders\hrule height \headrulewidth\hfill}}

% ----- Заголовки -----
\titleformat{\section}{\large\bfseries\color{accent}}{\thesection.}{0.6em}{}
\titleformat{\subsection}{\bfseries}{\thesubsection.}{0.5em}{}

% ----- Теоремные окружения (стиль конспекта) -----
\theoremstyle{definition}
\newtheorem{definition}{Определение}[section]
\newtheorem{example}[definition]{Пример}
\newtheorem{remark}[definition]{Замечание}
\theoremstyle{plain}
\newtheorem{theorem}[definition]{Теорема}
\newtheorem{lemma}[definition]{Лемма}
\newtheorem{corollary}[definition]{Следствие}
\newtheorem{property}[definition]{Свойство}
\newtheorem{criterion}[definition]{Критерий}
\renewcommand{\proofname}{Доказательство}

% ----- Операторы и сокращения -----
\DeclareMathOperator{\tg}{tg}
\DeclareMathOperator{\ctg}{ctg}
\DeclareMathOperator{\arctg}{arctg}
\DeclareMathOperator{\arcctg}{arcctg}
\DeclareMathOperator{\sh}{sh}
\DeclareMathOperator{\ch}{ch}
\DeclareMathOperator{\Th}{th}
\DeclareMathOperator{\cth}{cth}
\DeclareMathOperator{\Const}{const}
\DeclareMathOperator{\sign}{sign}
\DeclareMathOperator{\rang}{rang}
\DeclareMathOperator{\diam}{diam}
\DeclareMathOperator{\Span}{span}
\DeclareMathOperator{\Ker}{Ker}
\DeclareMathOperator{\Img}{Im}
\DeclareMathOperator{\tr}{tr}
\newcommand{\R}{\mathbb{R}}
\newcommand{\C}{\mathbb{C}}
\newcommand{\N}{\mathbb{N}}
\newcommand{\Z}{\mathbb{Z}}
\newcommand{\Q}{\mathbb{Q}}
\newcommand{\dx}{\,dx}
\newcommand{\dt}{\,dt}
\newcommand{\du}{\,du}
\newcommand{\eps}{\varepsilon}
\renewcommand{\Re}{\operatorname{Re}}
\renewcommand{\Im}{\operatorname{Im}}
\renewcommand{\le}{\leqslant}
\renewcommand{\ge}{\geqslant}
\newcommand{\scal}[2]{\left( #1,\,#2\right)}
\DeclarePairedDelimiter{\abs}{\lvert}{\rvert}
\DeclarePairedDelimiter{\norm}{\lVert}{\rVert}

\begin{document}
\begin{center}
  {\LARGE\bfseries \doctitle}\\[4pt]
  {\large\itshape \coursename}
\end{center}
\vspace{6pt}

\section{Преобразование Абеля (суммирование по частям)}

Тонкие признаки сходимости позволяют исследовать ряды $\sum a_k b_k$, к которым неприменимы ни
признак Лейбница, ни признаки сравнения. В их основе лежит \emph{преобразование Абеля} —
дискретный аналог интегрирования по частям.

\begin{lemma}[преобразование Абеля]
Пусть $B_n=\sum_{k=1}^n b_k$, $B_0=0$. Тогда
\[
  \sum_{k=1}^n a_k b_k=a_n B_n+\sum_{k=1}^{n-1}(a_k-a_{k+1})B_k .
\]
\end{lemma}

\begin{proof}
Подставим $b_k=B_k-B_{k-1}$ и сгруппируем по $B_k$:
\[
  \sum_{k=1}^n a_k(B_k-B_{k-1})
   =\sum_{k=1}^{n-1}B_k(a_k-a_{k+1})+a_n B_n .
\]
Здесь $a_k$ играет роль «множителя», $b_k$ — «дифференциала»; разность $a_k-a_{k+1}$ заменяет
производную, а суммирование $B_k$ — первообразную. Это полностью повторяет схему
$\int u\,dv=uv-\int v\,du$.
\end{proof}

\section{Признак Дирихле}

\begin{theorem}[Дирихле]
Ряд $\displaystyle\sum_{k=1}^\infty a_k b_k$ сходится, если:
\begin{enumerate}[label=\arabic*),leftmargin=*]
  \item $a_k$ монотонно стремится к нулю ($a_k\downarrow0$);
  \item частичные суммы $B_n=\sum_{k=1}^n b_k$ ограничены в совокупности: $\abs{B_n}\le M$.
\end{enumerate}
\end{theorem}

\begin{proof}
По преобразованию Абеля $\displaystyle S_n=a_n B_n+\sum_{k=1}^{n-1}(a_k-a_{k+1})B_k$.
\emph{Первое слагаемое} $a_n B_n\to0$, так как $a_n\to0$, а $\abs{B_n}\le M$.
\emph{Второе слагаемое} — частичная сумма абсолютно сходящегося ряда:
\[
  \sum_{k=1}^{\infty}\abs{a_k-a_{k+1}}\abs{B_k}\le M\sum_{k=1}^{\infty}(a_k-a_{k+1})
   =M\,a_1<\infty
\]
(модули сняты в силу монотонности, ряд телескопический). Значит, существует
$\displaystyle\lim_{n\to\infty}S_n$.
\end{proof}

\begin{example}[классический пример $\sum\frac{\sin n}{n^p}$]
Ряд $\displaystyle\sum_{n=1}^\infty\frac{\sin n}{n^p}$ при $p>0$ сходится. Берём
$a_n=\dfrac1{n^p}\downarrow0$ и $b_n=\sin n$; частичные суммы синусов ограничены:
\[
  B_n=\sum_{k=1}^n\sin k=\frac{\cos\frac12-\cos\bigl(n+\frac12\bigr)}{2\sin\frac12},\qquad
  \abs{B_n}\le\frac{1}{\sin\frac12}.
\]
По признаку Дирихле ряд сходится. При $0<p\le1$ сходимость условная, при $p>1$ — абсолютная.
\end{example}

\section{Признак Абеля}

\begin{theorem}[Абель]
Ряд $\displaystyle\sum_{k=1}^\infty a_k b_k$ сходится, если:
\begin{enumerate}[label=\arabic*),leftmargin=*]
  \item $a_k$ монотонна и ограничена;
  \item ряд $\sum b_k$ сходится.
\end{enumerate}
\end{theorem}

\begin{proof}
Монотонная ограниченная последовательность $a_k$ имеет конечный предел $a$. Частичные суммы $B_n$
сходятся (ряд $\sum b_k$ сходится), значит ограничены: $\abs{B_k}\le M$. Снова по преобразованию
Абеля $S_n=a_n B_n+\sum_{k=1}^{n-1}(a_k-a_{k+1})B_k$: первое слагаемое имеет предел $a\cdot B$
(где $B=\sum b_k$), второе сходится абсолютно, как в признаке Дирихле. Поэтому $\lim S_n$
существует.
\end{proof}

\begin{remark}[различие условий]
В обоих признаках $a_k$ монотонна. Различие — в требованиях:
\begin{itemize}[leftmargin=*]
  \item \textbf{Дирихле}: $a_k\to0$ (сильное условие на $a_k$), но от $\sum b_k$ требуется лишь
        \emph{ограниченность} частичных сумм (ряд может расходиться, как $\sum\sin k$);
  \item \textbf{Абель}: $a_k$ лишь \emph{ограничена} (предел может быть ненулевым), зато $\sum b_k$
        обязан \emph{сходиться}.
\end{itemize}
Признак Абеля сводится к Дирихле: если $a_k\to a$, то $a_k-a\to0$ монотонно, а
$\sum(a_k-a)b_k$ исследуется по Дирихле, $\sum a\,b_k$ сходится по условию.
\end{remark}

\begin{example}
$\displaystyle\sum_{k=1}^\infty\frac{(-1)^k}{\sqrt k}\Bigl(1+\frac1k\Bigr)^k$: множитель
$a_k=\bigl(1+\tfrac1k\bigr)^k$ возрастает и ограничен ($0\le a_k<e$), а ряд
$\sum\dfrac{(-1)^k}{\sqrt k}$ сходится по признаку Лейбница. По признаку Абеля исходный ряд
сходится.
\end{example}

\end{document}
